\section{Conclusion and Future Work}
\label{sec:conclusion}

We presented a rigorous study of geometric model merging, formulated through \textbf{RIMO} (\textbf{R}iemannian \textbf{I}sometry-respecting \textbf{M}anifold \textbf{O}perations), which reveals critical boundaries on non-linear manifolds. First, Riemannian merging is highly sensitive to parameter orthogonality: training with soft orthogonal regularization ($\lambda_{ortho} = 2.0$) minimizes Euclidean residuals, doubling the merged accuracy from $42.07\%$ to $84.55\%$. Second, tangent-space spectral balancing (RIMO with $t > 1.0$) catastrophically collapses accuracy to $13.66\%$ because the non-linear Cayley map translates tangent-space noise into spurious rotations across inactive dimensions.

Our findings demonstrate that geometric merging must respect the intrinsic rank and manifold constraints of task representations. To build upon these insights, we identify several highly promising pathways for future research:
\textbf{(1) Scaling to Modern Architectures and PEFT/LoRA:} Evaluating RIMO and RIMO-Pruned on modern ViTs and LLMs fine-tuned via LoRA is a natural and highly scalable extension.
\textbf{(2) Integration of Hard Orthogonal Constraints:} Exploring hard orthogonal constraints during training (e.g., Riemannian SGD on the Stiefel manifold or geotorch) to enforce $\|\rho_k\|_F = 0$ exactly will completely decouple residual effects from manifold interpolation.
\textbf{(3) Adaptive and Saliency-Based Tangent-Space Pruning:} Refining RIMO-Pruned by designing adaptive thresholds or incorporating parameter saliency (e.g., Fisher Information or Hessian curvature) to selectively preserve critical task directions.
\textbf{(4) Extension to Other Riemannian Lie Groups and Manifolds:} Generalizing RIMO to other non-Euclidean structures, such as the Unitary group $\mathrm{U}(d)$, hyperbolic spaces $\mathbb{H}^n$, the Special Euclidean group $\mathrm{SE}(d)$, or the Symplectic group $\mathrm{Sp}(2d)$ (which can preserve critical physical neural properties, such as energy conservation and stability in recurrent dynamics) to preserve physical and geometric neural structures.
